\vsssub
\subsubsection{~$S_{bot}$：\showex\ 底摩擦} \label{sec:BT4}
\vsssub

\opthead{BT4}{Crest model}{F. Ardhuin}

\noindent
针对砂质海底，一个更现实的参数化基于 \cite{art:GM79} 的涡黏模型，以及包含沙波形成
和向片流过渡的粗糙度参数化。\cite{tol:JPO94} 的参数化由 \cite{art:Aea03a}
根据北卡罗来纳大陆架上 DUCK'94 和 SHOWEX 实验的现场测量作了调整。该参数化也被
适配到 \ws\ 中，并加入了 \cite{tol:CE95} 给出的水深变率子网格参数化。该参数化
通过开关 BT4 启用。

源项可写为

%------------------------%
% SHOWEX bottom friction %
%------------------------%
% eq:SHOWEX_bot

\begin{equation}
\cS_{bot}(k,\theta) = - f_e u_b \: \frac{\sigma^2}{2 g \sinh^2(kd)} \: N(k,\theta)
\: , \label{eq:SHOWEX_bot}
\end{equation}

\noindent
其中 $f_e$ 为耗散因子，它是底部轨道位移均方根振幅 $a_b$ 和 Nikuradse 粗糙长度
$k_N$ 的函数；$u_b$ 为底部轨道速度的均方根值。

当前海床粗糙度参数化
(\ref{eq:SHOWEX_kr})--(\ref{eq:SHOWEX_krr}) 含有表 \ref{tab:BT4} 所列的
七个经验系数。

\begin{table} \begin{center}
\begin{tabular}{|l|c|c|c|c|c|} \hline \hline
参数         &  WWATCH 变量           & namelist &  SHOWEX  &  \cite{tol:JPO94} \\
\hline
  $A_1$ &  RIPFAC1                    & BT4 & 0.4     & 1.5    \\
  $A_2$ &  RIPFAC2                    & BT4 & -2.5    & -2.5   \\
  $A_3$ &  RIPFAC3                    & BT4 & 1.2     & 1.2     \\
  $A_4$ &  RIPFAC4                    & BT4 & 0.05    & 0.0     \\
  $\sigma_d$ &  SIGDEPTH              & BT4 & 0.05    & 用户定义  \\
  $A_5$ &  BOTROUGHMIN                & BT4 & 0.01    & 0.0     \\
  $A_6$ &  BOTROUGHFAC                & BT4 & 1.00    & 0.0     \\
\hline
\end{tabular} \end{center}
\caption{SHOWEX 底摩擦（默认值）以及 \cite{tol:JPO94} 所用原始参数值。
源项参数可通过 BT4 namelist 修改。请注意，变量名只适用于 namelist。在源项模块中，
这七个变量包含在数组 SBTCX 中。} \label{tab:BT4}
\botline
\end{table}

粗糙度 $k_{N}$ 分解为沙波粗糙度 $k_{r}$ 和片流粗糙度 $k_{s}$：
\begin{eqnarray}
k_{r} &=&a_{b}\times A_{1}\left( \frac{\psi }{\psi _{c}}\right) ^{A_{2}}, \label{eq:SHOWEX_kr}\\
k_{s} &=&0.57\frac{u_{b}^{2.8}}{\left[ g\left( s-1\right) \right] ^{1.4}}%
\frac{a_{b}^{-0.4}}{\left( 2\pi \right) ^{2}}\label{eq:SHOWEX_ks}.
\end{eqnarray}
在式 (\ref{eq:SHOWEX_kr}) 和 (\ref{eq:SHOWEX_ks}) 中，$A_1$ 和 $A_2$ 为经验常数，
$s$ 为泥沙比重，$\psi$ 为 Shields 数，由 $u_b$ 和中值砂粒直径 $D_{50}$ 确定：
\begin{equation}
\psi =f_{w}^{\prime }u_{b}^{2}/\left[g\left( s-1\right) D_{50}\right],
\end{equation}
其中 $f_{w}^{\prime }$ 为砂粒摩擦因子（其确定方式与 $f_e$ 相同，但以 $D_{50}$
而不是 $k_r$ 作为底部粗糙度）；$\psi _{c}$ 为平床上正弦波作用下泥沙开始运动的
临界 Shields 数。我们使用一个解析拟合式 \citep{bk:Soul97}：
\begin{eqnarray}
\psi _{c} &=&\frac{0.3}{1+1.2D_{*}}+0.055\left[ 1-\exp \left(
-0.02D_{*}\right) \right]\label{Soulsby_psic} , \\
D_{*} &=&D_{50}\left[ \frac{g\left( s-1\right) }{\nu ^{2}}\right] ^{1/3},
\end{eqnarray}
其中 $\nu $ 为水的运动黏性系数。

当波动不足以产生涡旋沙波时，即 Shields 数低于阈值 $\psi _{{\rm rr}}$ 时，
$k_{N}$ 由遗留沙波粗糙度 $k_{{\rm rr}}$ 给出。该阈值为
\begin{equation}
\psi _{{\rm rr}}=A_{3}\psi _{c}.
\end{equation}
低于该阈值时，$k_{N}$ 给定为
\begin{equation}
k_{{\rm rr}}=\max \left\{ A_5 {\rm m,} A_6 D_{50}, A_{4}a_{b}\right\}
{\rm for\ }\psi <\psi _{{\rm rr}}.\label{eq:SHOWEX_krr}
\end{equation}
